\documentclass[12pt,a4paper]{article}
\usepackage[T2A]{fontenc}
\usepackage[utf8]{inputenc}
\usepackage[russian]{babel}
\usepackage{float}
\usepackage{amsmath, amssymb}
\usepackage{geometry}
\usepackage{tikz}
\usepackage{tikz-3dplot}
\usepackage{caption}
\usetikzlibrary{arrows.meta, calc, decorations.pathmorphing}
\geometry{margin=2cm}

\newcommand{\vn}{\bar n}
\newcommand{\vq}{\bar q}
\newcommand{\vr}{\bar r}
\newcommand{\vi}{\bar i}
\newcommand{\vj}{\bar j}
\newcommand{\vk}{\bar k}

\begin{document}

\setcounter{page}{70}

\section*{ЛЕКЦИЯ 10. УРАВНЕНИЯ ПОВЕРХНОСТИ И ЛИНИИ В ПРОСТРАНСТВЕ}

\subsection*{10.1. Уравнения плоскости в пространстве}

\begin{figure}[H]
  \centering
  \begin{minipage}{0.4\textwidth}
    \centering
    \tdplotsetmaincoords{65}{115}
    \begin{tikzpicture}[scale=0.9, tdplot_main_coords, >=Latex]
      \draw[->] (0,0,0) -- (3,0,0) node[below left] {$x$};
      \draw[->] (0,0,0) -- (0,3,0) node[right] {$y$};
      \draw[->] (0,0,0) -- (0,0,3) node[above] {$z$};
      \filldraw[fill=gray!15, draw=black] (0.7,0.5,0.4) --
      (3.0,0.8,0.5) -- (2.5,2.4,1.5) -- (0.2,2.1,1.4) -- cycle;
      \fill (1.8,1.5,1.1) circle (1.2pt);
      \node[above] at (1.8,1.5,1.1) {$M(x,y,z)$};
      \node at (2.6,2.0,1.6) {$S$};
    \end{tikzpicture}
    \caption*{Рисунок 10.1}
  \end{minipage}%
  \begin{minipage}{0.55\textwidth}
    \textbf{Определение.} Уравнение
    $$
    \Phi(x,y,z)=0
    $$
    называется уравнением поверхности $S$, если этому уравнению
    удовлетворяют координаты любой точки этой поверхности, и не
    удовлетворяют координаты точек, не принадлежащих поверхности $S$.
  \end{minipage}
\end{figure}

\subsection*{\textbf{Уравнение плоскости в пространстве}}

\textbf{Теорема.} Любая плоскость в пространстве задается уравнением вида
\begin{equation}
  Ax+By+Cz+D=0. \tag{10.1}
\end{equation}
И обратно: всякое уравнение вида (10.1) определяет в пространстве плоскость.

\subsection*{\textbf{Доказательство}}

$\vn\perp Q$, $\vn$ --- нормальный вектор.

Возьмем произвольную точку
$$
M(x,y,z)\in Q.
$$

\begin{figure}[H]
  \centering
  \begin{minipage}{0.4\textwidth}
    \centering
    \tdplotsetmaincoords{65}{115}
    \begin{tikzpicture}[scale=0.85, tdplot_main_coords, >=Latex]
      \draw[->] (0,0,0) -- (3.5,0,0) node[below left] {$x$};
      \draw[->] (0,0,0) -- (0,3.5,0) node[right] {$y$};
      \draw[->] (0,0,0) -- (0,0,3.2) node[above] {$z$};
      \filldraw[fill=gray!15, draw=black] (0.5,0.4,0.6) --
      (3.2,0.5,0.7) -- (2.9,2.7,1.4) -- (0.2,2.4,1.3) -- cycle;
      \node at (2.9,2.25,1.45) {$Q$};
      \coordinate (M0) at (1.1,1.1,0.9);
      \coordinate (M) at (2.2,1.7,1.15);
      \fill (M0) circle (1.2pt);
      \fill (M) circle (1.2pt);
      \node[below] at (M0) {$M_0(x_0,y_0,z_0)$};
      \node[above] at (M) {$M(x,y,z)$};
      \draw[->, thick] (M0) -- ++(0.1,0.2,1.4) node[above] {$\vn(A,B,C)$};
    \end{tikzpicture}
    \caption*{Рисунок 10.2}
  \end{minipage}%
  \begin{minipage}{0.55\textwidth}
    $$
    \text{т. }M(x,y,z)\in Q
    \Leftrightarrow \vn\perp \overline{M_0M}
    \Leftrightarrow
    $$
    $$
    \Leftrightarrow (\vn,\overline{M_0M})=0
    \Leftrightarrow
    $$
    \begin{equation}
      \Leftrightarrow A(x-x_0)+B(y-y_0)+C(z-z_0)=0. \tag{10.2}
    \end{equation}
  \end{minipage}
\end{figure}

Раскроем скобки:
$$
Ax-Ax_0+By-By_0+Cz-Cz_0=
$$
$$
=[-Ax_0-By_0-Cz_0=D]=0\Rightarrow
$$
$$
\Rightarrow Ax+By+Cz+D=0.
$$

Проводя обратные рассуждения, получим, что уравнение вида (10.1)
определяет в пространстве плоскость.

\textbf{ч.т.д.}

(10.1) --- общее уравнение плоскости.

(10.2) --- уравнение плоскости, проходящей через точку $M_0(x_0,y_0,z_0)$
перпендикулярно вектору $\vn(A,B,C)$.

Из (10.1) $\Rightarrow$
$$
\frac{x}{-D/A}+\frac{y}{-D/B}+\frac{z}{-D/C}=1.
$$

Пусть
$$
a=-\frac DA;\qquad b=-\frac DB;\qquad c=-\frac DC,
$$
получим уравнение
\begin{equation}
  \frac xa+\frac yb+\frac zc=1. \tag{10.3}
\end{equation}

\begin{figure}[H]
  \centering
  \begin{minipage}{0.4\textwidth}
    \centering
    \tdplotsetmaincoords{65}{115}
    \begin{tikzpicture}[scale=0.85, tdplot_main_coords, >=Latex]
      \draw[->] (0,0,0) -- (3.7,0,0) node[below left] {$x$};
      \draw[->] (0,0,0) -- (0,3.7,0) node[right] {$y$};
      \draw[->] (0,0,0) -- (0,0,3.7) node[above] {$z$};
      \coordinate (A) at (3,0,0);
      \coordinate (B) at (0,3,0);
      \coordinate (C) at (0,0,3);
      \filldraw[fill=gray!15, draw=black] (A) -- (B) -- (C) -- cycle;
      \node[below] at (A) {$a$};
      \node[right] at (B) {$b$};
      \node[left] at (C) {$c$};
    \end{tikzpicture}
    \caption*{Рисунок 10.3}
  \end{minipage}%
  \begin{minipage}{0.55\textwidth}
    (10.3) --- уравнение плоскости в отрезках.
  \end{minipage}
\end{figure}

\subsection*{\textbf{Уравнение плоскости, проходящей через три точки}}

\begin{equation}
  \begin{vmatrix}
    x-x_1 & y-y_1 & z-z_1\\
    x_2-x_1 & y_2-y_1 & z_2-z_1\\
    x_3-x_1 & y_3-y_1 & z_3-z_1
  \end{vmatrix}=0. \tag{10.4}
\end{equation}

\begin{figure}[H]
  \centering
  \begin{minipage}{0.4\textwidth}
    \centering
    \tdplotsetmaincoords{65}{115}
    \begin{tikzpicture}[scale=0.85, tdplot_main_coords, >=Latex]
      \draw[->] (0,0,0) -- (3.3,0,0) node[below left] {$x$};
      \draw[->] (0,0,0) -- (0,3.3,0) node[right] {$y$};
      \draw[->] (0,0,0) -- (0,0,3.0) node[above] {$z$};
      \coordinate (A) at (0.7,0.8,0.8);
      \coordinate (B) at (3.0,0.8,1.1);
      \coordinate (C) at (1.1,2.6,1.5);
      \filldraw[fill=gray!15, draw=black] (A) -- (B) -- (C) -- cycle;
      \fill (A) circle (1.2pt) node[below] {$A$};
      \fill (B) circle (1.2pt) node[right] {$B$};
      \fill (C) circle (1.2pt) node[above] {$C$};
    \end{tikzpicture}
    \caption*{Рисунок 10.4}
  \end{minipage}%
  \begin{minipage}{0.55\textwidth}
    (10.4) --- уравнение плоскости $ABC$.
  \end{minipage}
\end{figure}

\subsection*{\textbf{Нормальное уравнение плоскости}}

\begin{equation}
  x\cos\alpha+y\cos\beta+z\cos\gamma-p=0. \tag{10.5}
\end{equation}

(10.5) --- нормальное уравнение плоскости, где $\cos\alpha$, $\cos\beta$,
$\cos\gamma$ --- направляющие косинусы нормального вектора,
направленного из начала координат в сторону плоскости, а $p$ ---
расстояние от начала координат до плоскости.

\subsection*{\textbf{Расстояние от точки до плоскости}}

$$
Q:Ax+By+Cz+D=0.
$$

\begin{equation}
  d=\left|\frac{Ax_0+By_0+Cz_0+D}{\sqrt{A^2+B^2+C^2}}\right|. \tag{10.5a}
\end{equation}

\begin{figure}[H]
  \centering
  \begin{minipage}{0.4\textwidth}
    \centering
    \tdplotsetmaincoords{65}{115}
    \begin{tikzpicture}[scale=0.85, tdplot_main_coords, >=Latex]
      \filldraw[fill=gray!15, draw=black] (0.2,0.5,0.4) --
      (3.4,0.7,0.5) -- (3.1,2.7,1.2) -- (0,2.4,1.1) -- cycle;
      \node at (3.1,2.25,1.25) {$Q$};
      \coordinate (P) at (1.8,1.5,0.85);
      \coordinate (M) at (1.8,1.5,2.7);
      \fill (M) circle (1.2pt) node[above] {$M_0(x_0,y_0,z_0)$};
      \draw[dashed] (M) -- (P) node[midway, right] {$d$};
      \fill (P) circle (1.2pt);
    \end{tikzpicture}
    \caption*{Рисунок 10.5}
  \end{minipage}%
  \begin{minipage}{0.55\textwidth}
    Расстояние от точки $M_0(x_0,y_0,z_0)$ до плоскости $Q$ выражается
    через подстановку координат точки в левую часть общего уравнения
    плоскости.
  \end{minipage}
\end{figure}

\subsection*{\textbf{Угол между плоскостями}}

$$
Q_1:A_1x+B_1y+C_1z+D_1=0;
$$
$$
\alpha=(\hat Q_1,Q_2);
$$
$$
Q_2:A_2x+B_2y+C_2z+D_2=0.
$$

$$
\alpha=\arccos\left|
\frac{A_1A_2+B_1B_2+C_1C_2}
{\sqrt{A_1^2+B_1^2+C_1^2}\cdot\sqrt{A_2^2+B_2^2+C_2^2}}
\right|.
$$

$$
A_1A_2+B_1B_2+C_1C_2=0
$$
--- условие перпендикулярности.

$$
\frac{A_1}{A_2}=\frac{B_1}{B_2}=\frac{C_1}{C_2}
$$
--- условие параллельности.

\subsection*{10.2. Уравнения прямой в пространстве}

\subsection*{\textbf{Уравнение линии в пространстве}}

\begin{equation}
  \begin{cases}
    x=x(t),\\
    y=y(t),\\
    z=z(t).
  \end{cases}
  \qquad t\in[\alpha,\beta]. \tag{10.6}
\end{equation}

\begin{figure}[H]
  \centering
  \begin{minipage}{0.4\textwidth}
    \centering
    \tdplotsetmaincoords{65}{115}
    \begin{tikzpicture}[scale=0.9, tdplot_main_coords, >=Latex]
      \draw[->] (0,0,0) -- (3.4,0,0) node[below left] {$x$};
      \draw[->] (0,0,0) -- (0,3.4,0) node[right] {$y$};
      \draw[->] (0,0,0) -- (0,0,3.2) node[above] {$z$};
      \draw[thick] (0.6,0.5,0.4) .. controls (1.4,2.4,0.9) and
      (2.2,0.9,2.1) .. (3.0,2.6,2.4) node[right] {$L$};
      \coordinate (M) at (2.2,1.25,1.9);
      \fill (M) circle (1.2pt);
      \node[above] at (M) {$M(x,y,z)$};
      \draw[->, thick] (0,0,0) -- (M) node[midway, below] {$\vr$};
      \node[below left] at (0,0,0) {$0$};
    \end{tikzpicture}
    \caption*{Рисунок 10.6}
  \end{minipage}%
  \begin{minipage}{0.55\textwidth}
    (10.6) --- параметрическое уравнение линии в пространстве.
  \end{minipage}
\end{figure}

\begin{equation}
  \vr(t)=x(t)\vi+y(t)\vj+z(t)\vk
  \quad \text{или}\quad
  \vr(t)=(x(t),y(t),z(t)). \tag{10.7}
\end{equation}

(10.7) --- векторное уравнение линии $L$.

Функция $\vr(t)$ называется вектор-функцией скалярного аргумента $t$.

\subsection*{\textbf{Уравнение прямой в пространстве}}

\begin{equation}
  \frac{x-x_0}{l}=\frac{y-y_0}{m}=\frac{z-z_0}{n}. \tag{10.8}
\end{equation}

\begin{figure}[H]
  \centering
  \begin{minipage}{0.4\textwidth}
    \centering
    \tdplotsetmaincoords{65}{115}
    \begin{tikzpicture}[scale=0.9, tdplot_main_coords, >=Latex]
      \draw[->] (0,0,0) -- (3.4,0,0) node[below left] {$x$};
      \draw[->] (0,0,0) -- (0,3.4,0) node[right] {$y$};
      \draw[->] (0,0,0) -- (0,0,3.2) node[above] {$z$};
      \draw[thick] (0.5,0.7,0.4) -- (3.0,2.5,2.1) node[right] {$L$};
      \draw[->, thick] (1.4,1.35,1.05) -- (2.3,2.0,1.65) node[above]
      {$\vq(l,m,n)$};
      \node[below left] at (0,0,0) {$0$};
    \end{tikzpicture}
    \caption*{Рисунок 10.7}
  \end{minipage}%
  \begin{minipage}{0.55\textwidth}
    (10.8) --- каноническое уравнение прямой, проходящей через точку
    $M_0(x_0,y_0,z_0)$ параллельно направляющему вектору
    $\vq(l,m,n)$.
  \end{minipage}
\end{figure}

\begin{equation}
  \begin{cases}
    x=x_0+l\cdot t,\\
    y=y_0+m\cdot t,\\
    z=z_0+n\cdot t.
  \end{cases}
  \tag{10.9}
\end{equation}

\begin{equation}
  \begin{cases}
    A_1x+B_1y+C_1z+D_1=0,\\
    A_2x+B_2y+C_2z+D_2=0.
  \end{cases}
  \tag{10.10}
\end{equation}

(10.10) --- уравнение прямой, как результат пересечения двух плоскостей.

(10.9) --- параметрическое уравнение прямой.

\subsection*{\textbf{Взаимное расположение двух прямых в пространстве}}

$$
L_1:\frac{x-x_1}{l_1}=\frac{y-y_1}{m_1}=\frac{z-z_1}{n_1};
$$
$$
L_2:\frac{x-x_2}{l_2}=\frac{y-y_2}{m_2}=\frac{z-z_2}{n_2}.
$$

1) $L_1\parallel L_2\Leftrightarrow$
$$
\frac{l_1}{l_2}=\frac{m_1}{m_2}=\frac{n_1}{n_2};
$$

2) $L_1\perp L_2\Leftrightarrow l_1l_2+m_1m_2+n_1n_2=0$;

3) $\alpha=(\hat L_1,L_2)$
$$
\alpha=\arccos\left|
\frac{l_1l_2+m_1m_2+n_1n_2}
{\sqrt{l_1^2+m_1^2+n_1^2}\cdot\sqrt{l_2^2+m_2^2+n_2^2}}
\right|;
$$

4) Если $L_1\nparallel L_2$, то $L_1$ и $L_2$ --- пересекаются $\Leftrightarrow$
$$
\Leftrightarrow
\Delta=
\begin{vmatrix}
x_2-x_1 & y_2-y_1 & z_2-z_1\\
l_1 & m_1 & n_1\\
l_2 & m_2 & n_2
\end{vmatrix}=0;
$$

5) $L_1$ и $L_2$ --- скрещиваются $\Leftrightarrow \Delta\ne0$.

$$
Q:Ax+By+Cz+D=0;
\qquad
L:\frac{x-x_0}{l}=\frac{y-y_0}{m}=\frac{z-z_0}{n};
$$

$$
L\parallel Q\Leftrightarrow Al+Bm+Cn=0,\qquad \vq\perp\vn;
$$

$$
L\perp Q\Leftrightarrow \frac Al=\frac Bm=\frac Cn,\qquad \vq\parallel\vn.
$$

$$
\varphi=(\hat\vn,\vq),
\qquad
\alpha=(\hat L,Q).
$$

\begin{figure}[H]
\centering
\begin{minipage}{0.4\textwidth}
\centering
\tdplotsetmaincoords{65}{115}
\begin{tikzpicture}[scale=0.9, tdplot_main_coords, >=Latex]
\filldraw[fill=gray!15, draw=black] (0.2,0.5,0.3) -- (3.6,0.7,0.4) --
(3.2,2.6,1.1) -- (0,2.4,1.0) -- cycle;
\node at (3.1,2.3,1.2) {$Q$};
\draw[thick] (0.5,0.7,1.8) -- (3.3,2.5,0.1) node[right] {$L$};
\draw[->, thick] (1.6,1.4,0.75) -- (2.3,1.85,0.35) node[above] {$\vq$};
\draw[->, thick] (1.6,1.4,0.75) -- (1.7,1.55,2.3) node[above] {$\vn$};
\node at (1.9,1.65,1.5) {$\varphi$};
\node at (2.4,2.0,0.65) {$\alpha$};
\end{tikzpicture}
\caption*{Рисунок 10.8}
\end{minipage}%
\begin{minipage}{0.55\textwidth}
$$
\cos\varphi=\frac{(\vn,\vq)}{|\vn||\vq|}
=\cos\left(\frac\pi2-\alpha\right)=\sin\alpha
\Rightarrow
$$
$$
\Rightarrow
\alpha=\arcsin\left|
\frac{Al+Bm+Cn}
{\sqrt{A^2+B^2+C^2}\cdot\sqrt{l^2+m^2+n^2}}
\right|.
$$
\end{minipage}
\end{figure}

\subsection*{\textbf{Пример}}

Составить уравнение плоскости $ABC$, если
$$
A(1;1;-1);\qquad B(2;0;1);\qquad C(-1;2;3).
$$

\subsection*{\textbf{Решение}}

$$
\overline{AB}(1;-1;2),\qquad \overline{AC}(-2;1;4);
$$

$$
[\overline{AB};\overline{AC}]=
\begin{vmatrix}
\vi & \vj & \vk\\
1 & -1 & 2\\
-2 & 1 & 4
\end{vmatrix}
=
\vi
\begin{vmatrix}-1&2\\1&4
\end{vmatrix}
-\vj
\begin{vmatrix}1&2\\-2&4
\end{vmatrix}
+\vk
\begin{vmatrix}1&-1\\-2&1
\end{vmatrix}
$$
$$
=-6\vi-8\vj-\vk\Rightarrow \vn(6;8;1).
$$

$$
6(x-1)+8(y-1)+1(z+1)=
$$
$$
=6x-6+8y-8+z+1=0.
$$

$$
6x+8y+z-13=0
$$
--- уравнение плоскости $ABC$.

\subsection*{\textbf{Пример}}

Найти угол между плоскостями $x-y+2z-11=0$ и $2x+y+1=0$.

\subsection*{\textbf{Решение}}

$$
\vn_1(1;-1;2),\qquad \vn_2(2;1;0).
$$

$$
\cos\alpha=
\frac{(\vn_1,\vn_2)}{|\vn_1||\vn_2|}
=
\left|
\frac{1\cdot2-1\cdot1+2\cdot0}{\sqrt{1+1+4}\cdot\sqrt{4+1}}
\right|
=
\frac1{\sqrt6\cdot\sqrt5}
=
\frac1{\sqrt{30}}.
$$

$$
\Rightarrow \alpha=\arccos\frac1{\sqrt{30}}.
$$

\subsection*{\textbf{Пример}}

Найти расстояние между плоскостями
$$
Q_1:3x-y+2z-5=0
$$
и
$$
Q_2:-6x+2y-4z+7=0.
$$

\subsection*{\textbf{Решение}}

Найдем произвольную точку $M\in Q_1$: $M(1;0;1)$. Найдем расстояние от
точки $M$ до плоскости $Q_2$:
$$
d=\left|\frac{-6\cdot1+2\cdot0-4\cdot1+7}{\sqrt{36+4+16}}\right|
=
\left|\frac{-6-4+7}{2\sqrt{9+1+4}}\right|
=
\frac3{2\sqrt{14}}.
$$

\subsection*{\textbf{Пример}}

Найти каноническое и параметрическое уравнения прямой $AB$, если
$$
A(3;2;1)\quad \text{и}\quad B(1;-1;2).
$$

\subsection*{\textbf{Решение}}

$$
\overline{AB}(-2;-3;1)\Rightarrow \vq(2;3;-1).
$$

$$
\frac{x-1}{2}=\frac{y+1}{3}=\frac{z-2}{-1}
$$
--- каноническое уравнение.

$$
\frac{x-1}{2}=t\Rightarrow x=2t+1;
$$
$$
\frac{y+1}{3}=t\Rightarrow y=3t-1;
$$
$$
\frac{z-2}{-1}=t\Rightarrow z=-t+2;
$$

$$
\begin{cases}
x=1+2t,\\
y=-1+3t,\\
z=2-t.
\end{cases}
$$
--- параметрическое уравнение.

\subsection*{\textbf{Пример}}

Найти каноническое и параметрическое уравнения прямой
$$
\begin{cases}
x+2y-3z-2=0,\\
x+y-2z-6=0.
\end{cases}
$$

\subsection*{\textbf{Решение}}

$$
\vn_1(1;2;-3),\qquad \vn_2(1;1;-2).
$$

$$
[\vn_1;\vn_2]=
\begin{vmatrix}
\vi & \vj & \vk\\
1 & 2 & -3\\
1 & 1 & -2
\end{vmatrix}
=
\vi
\begin{vmatrix}2&-3\\1&-2
\end{vmatrix}
-\vj
\begin{vmatrix}1&-3\\1&-2
\end{vmatrix}
+\vk
\begin{vmatrix}1&2\\1&1
\end{vmatrix}
$$
$$
=-\vi-\vj-\vk\Rightarrow \vq(1;1;1)
$$
--- направляющий вектор прямой.

Найдем точку $M$, лежащую на прямой. Пусть $y=0\Rightarrow$
$$
\Rightarrow
\begin{cases}
x-3z-2=0,\quad (1)\\
x-2z-6=0.\quad (2)
\end{cases}
$$

$$(2)-(1): z-4=0\Rightarrow z=4;$$
$$
x=3z+2=14\Rightarrow M(14;0;4).
$$

Так как $\vq(1;1;1)$ и $M(14;0;4)$:
$$
\frac{x-14}{1}=\frac{y}{1}=\frac{z-4}{1}
$$
--- каноническое уравнение прямой.

$$
\begin{cases}
x=14+t,\\
y=t,\\
z=4+t.
\end{cases}
$$
--- параметрическое уравнение прямой.

\subsection*{\textbf{Пример}}

Найти проекцию точки $A(6;-9;8)$ на плоскость $Q$, заданную уравнением
$$
x-2y+z-2=0.
$$

\subsection*{\textbf{Решение}}

Пусть точка $B$ --- проекция точки $A$ на плоскость $Q$.

Тогда прямая $AB\perp Q$.

Найдем уравнение прямой $AB$:
$$
\frac{x-6}{1}=\frac{y+9}{-2}=\frac{z-8}{1}
$$
--- каноническое уравнение прямой $AB$.

Найдем координаты точки $B$ --- точки пересечения прямой $AB$ с плоскостью
$Q$.

Для этого сначала запишем параметрическое уравнение прямой $AB$.
$$
\frac{x-6}{1}=\frac{y+9}{-2}=\frac{z-8}{1}.
$$

$$
\Rightarrow
\begin{cases}
x=6+t,\\
y=-9-2t,\\
z=8+t.
\end{cases}
$$
--- параметрическое уравнение прямой $AB$.

Далее подставим выражения для $x$, $y$ и $z$ в уравнение плоскости $Q$.
$$
Q:x-2y+z-2=0;
\qquad
AB:
\begin{cases}
x=6+t,\\
y=-9-2t,\\
z=8+t;
\end{cases}
$$

$$
6+t-2(-9-2t)+8+t-2=0;
$$
$$
6+t+18+4t+8+t-2=0;
$$
$$
6t+30=0\Rightarrow t=-5\Rightarrow B(1;1;3).
$$

\subsection*{\textbf{Пример}}

Найти угол между прямой $L$ и плоскостью $Q$, если
$$
L:\frac{x+2}{2}=\frac{y-2}{-1}=\frac{z}{3};
\qquad
Q:3x+y-2z-12=0.
$$

\subsection*{\textbf{Решение}}

$$
\sin\alpha=\left|\frac{(\vn,\vq)}{|\vn||\vq|}\right|,
$$
где $\vn(3;1;-2)$; $\vq(2;-1;3)$.

$$
\sin\alpha=
\left|
\frac{3\cdot2+1\cdot(-1)-2\cdot3}
{\sqrt{9+1+4}\sqrt{4+1+9}}
\right|
=
\left|\frac1{\sqrt{14}\sqrt{14}}\right|
=
\frac1{14};
$$

$$
\alpha=\arcsin\frac1{14}.
$$

\subsection*{\textbf{Пример}}

Найти проекцию точки $A(-5;-1;3)$ на прямую $L$, заданную уравнением
$$
\frac{x}{3}=\frac{y+1}{2}=\frac{z-2}{1}.
$$

\subsection*{\textbf{Решение}}

Найдем уравнение плоскости $Q$, проходящей через точку $A$ перпендикулярно
прямой $L$.

$$
3(x+5)+2(y+1)+z-3=0;
$$
$$
3x+2y+z+14=0.
$$

Пусть точка $B$ --- точка пересечения плоскости $Q$ с прямой $L$. Найдем эту
точку.

$$
\frac{x}{3}=\frac{y+1}{2}=\frac{z-2}{1}
\Rightarrow
\begin{cases}
x=3t,\\
y=-1+2t,\\
z=2+t;
\end{cases}
$$

$$
3x+2y+z+14=0;
$$
$$
3\cdot3t+2(-1+2t)+2+t+14=0;
$$
$$
9t-2+4t+2+t+14=0;
$$
$$
14t+14=0\Rightarrow t=-1\Rightarrow B(-3;-3;1).
$$

Точка $B$ --- проекция точки $A$ на прямую $L$.

\end{document}
